% ============================================================
%  Słownik formalizmu TGP — kanoniczne przejścia między formulacjami
%  TGP v1 — Appendix (do włączenia w main.tex)
% ============================================================
\section{Słownik formalizmu TGP}\label{app:formulation-dictionary}

\begin{quote}
\textbf{Cel}: TGP operuje kilkoma równoważnymi zestawami zmiennych
i~notacji. Niniejszy słownik definiuje kanoniczne przejścia,
wskazuje zakres stosowania każdej formulacji i~eliminuje
wieloznaczności dla recenzenta.
\end{quote}

% -------------------------------------------------------------------
\subsection{Formulacje i~ich domeny}
% -------------------------------------------------------------------

\begin{center}\small
\begin{tabular}{@{}lcccp{4.0cm}@{}}
\toprule
& \textbf{Form.\ A (kan.)} & \textbf{Form.\ Sub} & \textbf{Form.\ B (log)} & \textbf{Uwagi} \\
\midrule
\textbf{Zmienna pola} & $\varphi = \Phi/\PhiZero$ & $g = \Phi/\PhiZero$ & $g = \Phi/\PhiZero$ & Ta sama zmienna \\[4pt]
\textbf{Kinetyczne $K$} & $K(\varphi) = K_{\mathrm{geo}}\,\varphi^4$ & $K_{\mathrm{sub}}(g) = K_{\mathrm{geo}}\,g^2$ & $f(g) = 1 + 4\ln g$ & Trzy postaci \\[4pt]
\textbf{Eksponent $\alpha$} & $\alpha = 2$ & $\alpha = 1$ & $\alpha = 2$ (w~limicie) & $f = K'/K \cdot g/2$ \\[4pt]
\textbf{Ghost point} & $g^* = e^{-1/4} \approx 0.779$ & \textbf{brak} ($K > 0\;\forall g > 0$) & $g^* = e^{-1/4}$ & Krytyczna różnica \\[4pt]
\textbf{Relacja} & $K = K_{\mathrm{geo}}\,g^4$ & $K_{\mathrm{sub}} = K_{\mathrm{geo}}\,g^2$ & $f \approx 1 + 4\ln g$ & $f = K'/K\cdot g/2$ \\
\midrule
\textbf{ODE solitonu} &
  $g^4[g''+(2/r)g'] + 2g^3(g')^2 = V'$ &
  $g^2 g'' + g(g')^2 + (2/r)g^2 g' = V'$ &
  $f(g)[g''+(2/r)g'] + (2/g)(g')^2 = V'$ &
  \\[4pt]
\textbf{$r_{21}$ (phi-FP)} & 206.77 \checkmark & 206.77 \checkmark & 206.77 \checkmark & \textbf{Identyczny} \\[4pt]
\textbf{$r_{31}$ (tau)} & bariera: max$\sim$1571 & \textbf{3477} \checkmark & bariera & Sub jedyna pełna \\[4pt]
\textbf{$g_0^e$} & 0.834 & \textbf{0.869} & 0.899 & Zależy od formy \\
\midrule
\textbf{PPN $\gamma,\beta$} & $= 1$ & $= 1$ & $= 1$ & \textbf{Niezależne od $\alpha$} \\[4pt]
\textbf{$\kappa$} & $3/(4\PhiZero)$ & $3/(4\PhiZero)$ & $3/(4\PhiZero)$ & \textbf{Niezależne} \\[4pt]
\textbf{$n_s$, $r$} & 0.9662, 0.0033 & 0.9662, 0.0033 & 0.9662, 0.0033 & \textbf{Niezależne} \\[4pt]
\textbf{Koide $K$} & $2/3$ & $2/3$ & $2/3$ & \textbf{Niezależne} \\[4pt]
\textbf{$\alpha_s(M_Z)$} & 0.1184 & 0.1184 & 0.1184 & \textbf{Niezależne} \\
\bottomrule
\end{tabular}
\end{center}

\begin{remark}[Wybór kanonicznej formulacji]\label{rem:canonical-choice}
\textbf{Formulacja Substratowa} ($K_{\mathrm{sub}} = g^2$) jest
\textbf{preferowana} z~trzech powodów:
\begin{enumerate}[(i)]
  \item jest \textbf{ghost-free}: $K_{\mathrm{sub}}(g) = K_{\mathrm{geo}}\,g^2 > 0$
    dla wszystkich $g > 0$;
  \item reprodukuje \textbf{wszystkie trzy} masy leptonów
    ($r_{21}$ i~$r_{31}$) — jako jedyna;
  \item wynika bezpośrednio z~hamiltonianu substratu
    (lemat~\ref{lem:K_phi2}).
\end{enumerate}
Formulacje~A i~B są przybliżeniami lub historycznymi wariantami.
Formulacja~A jest poprawna w~sektorze słabego pola
($\varphi \sim 1$), ale niedziałająca w~sektorze solitonowym
($g_0 > g^*$).
\end{remark}

% -------------------------------------------------------------------
\subsection{Przejścia kanoniczne}
% -------------------------------------------------------------------

\begin{proposition}[Przejście Form.~Sub $\to$ Form.~A]\label{prop:sub-to-A}
Formulacja~A jest ,,modulacją'' Formulacji~Sub:
\begin{equation}
  K_A(\varphi) = \frac{d(g^2)}{dg} \cdot g \cdot K_{\mathrm{sub}}(g)
  = 2g \cdot g \cdot g^2 = 2g^4
  \quad \Rightarrow \quad
  K_A \propto g^4.
\end{equation}
Różnica wynika z~tego, czy kinetyczny grad pochodzi
z~$\phi_i$ (Sub) czy z~$\phi_i^2$ (A):
\begin{itemize}[nosep]
  \item Sub: $\mathcal{L} \supset J\sum \phi_i \phi_j
    \to K_{\mathrm{sub}} = g^2$
  \item A: $\mathcal{L} \supset K_{ij}\sum (\phi_i - \phi_j)^2
    \to K_A = g^4$ (z modulacją $K_{ij}$)
\end{itemize}
\end{proposition}

\begin{proposition}[Przejście Form.~A $\to$ Form.~B]\label{prop:A-to-B}
Formulacja~B jest rozwinięciem Taylora Formulacji~A wokół~$g = 1$:
\begin{equation}
  f(g) = \frac{d\ln K_A}{d\ln g}
  = \frac{g \cdot 4g^3}{g^4} = 4.
\end{equation}
Ogólniej, definiując $f(g) \equiv 1 + 2\alpha\ln g$
i~identyfikując z~$K_A(g)/K_A(1) \approx e^{2\alpha\ln g}$,
dostajemy $f(g) = 1 + 4\ln g$ ($\alpha = 2$) jako
pierwszorzędowe przybliżenie.
\end{proposition}

% -------------------------------------------------------------------
\subsection{Kiedy której formulacji użyć}
% -------------------------------------------------------------------

\begin{center}\small
\begin{tabular}{@{}lll@{}}
\toprule
\textbf{Sektor} & \textbf{Zalecana formulacja} & \textbf{Powód} \\
\midrule
Solitony (masy) & Sub & Ghost-free, reprodukuje $\tau$ \\
Grawitacja (PPN) & dowolna & Wynik niezależny od $\alpha$ \\
Kosmologia (FRW) & A lub Sub & Identyczne wyniki \\
Metryka & A & Historycznie, $\sqrt{-g} = c_0\varphi$ \\
Sektor cechowania & A lub Sub & Faza $\theta$ niezależna od modułu \\
\bottomrule
\end{tabular}
\end{center}

% -------------------------------------------------------------------
\subsection{Parametry: $\PhiZero$ i~$\PhiEff$}
% -------------------------------------------------------------------

\begin{remark}[Rozróżnienie $\PhiZero$ i~$\Phi_{\mathrm{eff}}$]
\begin{itemize}[nosep]
  \item $\PhiZero = 25$: wartość próżniowa pola $\Phi$
    (z~warunku $U'(\varphi_0) = 0$, $\varphi_0 = 1$,
    $\Phi = \PhiZero\varphi$). Jest parametrem \textbf{Warstwy~I}
    (strukturalnym) jeśli $\PhiZero = N_f^2 = (2N_c - 1)^2$.
  \item $\Phi_{\mathrm{eff}}$: ewentualna korekcja radiacyjna
    lub kosmologiczna do $\PhiZero$. W~bieżącej wersji TGP
    $\Phi_{\mathrm{eff}} = \PhiZero$ (brak poprawek).
  \item Hipoteza $a_\Gamma \cdot \PhiZero = 1$ łączy parametr
    substratu z~próżnią: weryfikowana z~danymi DESI
    (odchylenie $\sim 0.5\%$ na 2026-04).
\end{itemize}
\end{remark}

% -------------------------------------------------------------------
\begin{remark}[Update post-UV.3 (2026-05-02): $\Phi_0^{\text{bare}}$ vs $\PhiEff$]%
\label{rem:slownik-uv3-update}
% -------------------------------------------------------------------
Po zamknięciu UV.3 (\texttt{research/op-uv3-phi0-renormalization/}) rozróżnienie zostało wyostrzone:
\begin{itemize}[nosep]
  \item $\Phi_0^{\text{bare}} \equiv \Phi_0^{\text{bare}} = 168\cdot\Omega_\Lambda \approx 115$
        --- UV anchor (prze-renormalizowana wartość substratowa).
  \item $Z_\Phi = V(1)/P(1) = 14/3$ --- współczynnik renormalizacji
        wave-function (sympy LOCK; \texttt{op-uv3-phi0-renormalization/Phase2}).
  \item $\PhiEff = \Phi_0^{\text{bare}}/Z_\Phi \approx 24.65$ --- IR observable
        (poprzednio nazywane \emph{$\PhiZero$}). Numerycznie zbliżone
        do $\PhiZero=25$, ale strukturalnie distinct: $\PhiEff$ jest
        \emph{derived} z UV anchor + RG, nie input.
  \item Po γ.1 closure (\texttt{op-gamma1-phi-eff-anchor-resolution/}):
        $\PhiEff^{\text{pure}} = 8\pi \approx 25.13$ (czysta predykcja
        z $\Omega_\Lambda^{\text{pure}} = 2\pi/9$); po δ.1+δ.2 correction
        ($N_f = 5$ active QCD flavors): $\PhiEff^{\text{corr}} = (10/3)\cdot e^2 \approx 24.6302$,
        $\Omega_\Lambda^{\text{corr}} = 5e^2/54 \approx 0.68417$ (Planck 0.6847, drift $-0.07\sigma$).
\end{itemize}
\textbf{Zachowane:} $\PhiZero = 25$ jako \emph{shorthand} dla
$\PhiEff^{\text{corr}} \approx 24.63$ w pre-UV.3 derivations.
\textbf{Nowe:} cytaty post-2026-05-02 muszą rozróżniać $\Phi_0^{\text{bare}}$ vs $\PhiEff$
(zob. \texttt{core/sek00\_summary.tex} lin. 80--110, \texttt{rem:Lambda-post-uv3} w sek05).
\end{remark}

% -------------------------------------------------------------------
\subsection{Podsumowanie relacji kluczowych}
% -------------------------------------------------------------------

\begin{equation}\label{eq:dictionary-key}
\boxed{
\begin{aligned}
  K_{\mathrm{sub}}(g) &= K_{\mathrm{geo}}\,g^2
  & &\text{[substrat: } H = -J\sum \phi_i\phi_j]\\
  K_A(\varphi) &= K_{\mathrm{geo}}\,\varphi^4
  & &\text{[kanoniczny: modulacja } K_{ij}]\\
  f(g) &= 1 + 4\ln g
  & &\text{[logarytmiczny: Taylor wokół } g=1]\\[6pt]
  \alpha_{\mathrm{Sub}} &= 1, \quad \alpha_A = 2
  & &\text{[PPN/kosmo: niezależne od $\alpha$]}\\
  \kappa &= \tfrac{3}{4\PhiZero}
  & &\text{[z zunifikowanej akcji]}\\
  g_0^e &= 0.869\;\text{(Sub)},\; 0.834\;\text{(A)}
  & &\text{[bezwzgl. zależy od formy ODE]}\\
  r_{21} &= (A_\mu/A_e)^4 = 206.77
  & &\text{[niezależne od formy ODE — $\varphi$-FP]}
\end{aligned}
}
\end{equation}
